\section{\ENG{Failure Handling} \& \ENG{Resilience}}
\label{sec:failure_resilience}

Η ανθεκτικότητα της πλατφόρμας δεν οικοδομείται γύρω από σύνθετους \ENG{distributed failover} μηχανισμούς, αλλά γύρω από σαφή αποθήκευση της κατάστασης σε εξωτερικά αποθετήρια και επαναληπτικούς ελέγχους λειτουργίας. Επειδή τα διάφορα στάδια πληροφορίας αποθηκεύονται σε διαφορετικά συτήματα, η απώλεια ενός μεμονωμένου \ENG{process} δεν ταυτίζεται με απώλεια όλης της κατάστασης. Η αρχιτεκτονική αυτή δεν καταργεί τις αστοχίες, αλλά περιορίζει τη ζημία τους.

Σε επίπεδο κύκλου ζωής εκτέλεσης, ο σημαντικότερος μηχανισμός περιορισμού απώλειας είναι ο συνδυασμός περιοδικών \ENG{commits}, \ENG{sendBeacon} στο \ENG{beforeunload} και συγχρονισμού προς μόνιμη αποθήκευση τόσο στον κανονικό τερματισμό όσο και σε χρονικά παρωχημένες συνεδρίες. Ο \texttt{\ENG{RuntimeFlushScheduler}} εντοπίζει παρωχημένα ενεργά \ENG{launches} και τα προωθεί στον \texttt{\ENG{RuntimeFlushService}}, ο οποίος τα κλείνει ως \ENG{ABANDONED}. Επομένως, η πλατφόρμα έχει ένα σαφές σχέδιο αντιμετώπισης όχι μόνο των ομαλών τερματισμών αλλά και των σιωπηρών εγκαταλείψεων.

Στην υποδομή, η αξιοπιστία ενισχύεται από αυτοματοποιημένους ελέγχους περιβάλλοντος και \ENG{readiness}. Το \texttt{\ENG{validate-env.sh}} διασφαλίζει ότι το αρχείο \texttt{\ENG{.env}} υπάρχει, είναι μορφολογικά ορθό και περιέχει τις απαραίτητες μεταβλητές. Το \texttt{\ENG{smoke-check.sh}} ελέγχει σταδιακά αν τα \ENG{containers} εκκινούν σωστά και αν οι κρίσιμες υπηρεσίες είναι πραγματικά προσβάσιμες από το εσωτερικό δίκτυο \ENG{docker}. Το στοιχείο αυτό έχει ειδική σημασία, διότι μετατοπίζει την έννοια της αξιοπιστίας από θεωρητική επιθυμία σε επαναλήψιμη επιχειρησιακή διαδικασία.

Η παρούσα υλοποίηση, ωστόσο, δεν πρέπει να παρουσιαστεί ως πλήρες πλαίσιο ανθεκτικότητας. Δεν τεκμηριώνονται \ENG{circuit breakers}, ενεργοί \ENG{message replays}, \ENG{distributed sagas} ή πολύπλοκες στρατηγικές ανακατεύθυνσης αιτήσεων σε πολλαπλά \ENG{runtime instances}. Η ορθή περιγραφή είναι ότι η πλατφόρμα διαθέτει συγκεκριμένους, πρακτικούς μηχανισμούς περιορισμού απώλειας κατάστασης και ελέγχου διαθεσιμότητας, οι οποίοι εναρμονίζονται με το σημερινό επίπεδο πολυπλοκότητας του συστήματος.

Το Σχήμα~\ref{fig:failure_resilience} αποτυπώνει τους υλοποιημένους μηχανισμούς περιορισμού αστοχίας γύρω από \ENG{commit}, συγχρονισμό προς μόνιμη αποθήκευση και \ENG{health checking}. Ο Πίνακας~\ref{tab:operational_controls_and_checks} συγκεντρώνει τα βασικά λειτουργικά ελέγξιμα σημεία που στηρίζουν αυτή την επιχειρησιακή συμπεριφορά.

\begin{figure}[H]
  \centering
\begin{chapterfigurebox}
  \begin{tikzpicture}[
      font=\footnotesize,
      node distance=1.8cm and 2.5cm,
      box/.style={draw, rounded corners, fill=gray!5, align=center, text width=3.3cm, minimum height=1.05cm},
      emph/.style={draw, rounded corners, fill=gray!12, align=center, text width=3.6cm, minimum height=1.05cm, font=\footnotesize\bfseries},
      store/.style={draw, cylinder, shape border rotate=90, aspect=0.28, minimum height=1.2cm, minimum width=3.2cm, align=center, fill=gray!5},
      arrow/.style={-Latex, thick}
    ]
    \node[box] (browser) {\ENG{Browser / player}};
    \node[emph, right=of browser] (engine) {\ENG{Engine runtime}};
    \node[store, right=of engine] (redis) {\ENG{Redis}\\ενεργή κατάσταση};

    \node[box, below=1.9cm of browser] (scheduler) {\texttt{\ENG{RuntimeFlushScheduler}}};
    \node[emph, below=1.9cm of engine] (flush) {\texttt{\ENG{RuntimeFlushService}}};
    \node[store, below=1.9cm of redis] (postgres) {\ENG{Postgres}\\\ENG{attempts} + στιγμιότυπα};

    \node[box, below=1.9cm of flush] (checks) {\texttt{\ENG{validate-env.sh}}\\+\texttt{\ENG{smoke-check.sh}}};

    \draw[arrow] (browser) -- node[above, align=center]{περιοδικά \ENG{commits}\\ή \ENG{terminate}} (engine);
    \draw[arrow] (engine) -- (redis);
    \draw[arrow] (scheduler) -- node[above, align=center]{παρωχημένα\\\ENG{launches}} (flush);
    \draw[arrow] (engine) -- (flush);
    \draw[arrow] (flush) -- (postgres);
    \draw[arrow] (checks) -- (engine);
  \end{tikzpicture}
\end{chapterfigurebox}
  \caption{Υλοποιημένοι μηχανισμοί περιορισμού απώλειας κατάστασης και λειτουργικού ελέγχου γύρω από \ENG{commit}, συγχρονισμό προς μόνιμη αποθήκευση και ελέγχους ετοιμότητας.}
  \label{fig:failure_resilience}
\end{figure}

